\frontmatter

% ============================================================
% Halaman Sampul
% ============================================================
\begin{titlepage}
  \centering
  \vspace*{1cm}
  {\LARGE\bfseries Pemrograman Web\par}
  \vspace{0.5cm}
  {\large Buku Ajar Berbasis Outcome-Based Education (OBE)\par}
  \vspace{1cm}
  {\large Mata Kuliah: Pemrograman Web\par}
  {\large 2 SKS \textbar{} 16 Pertemuan\par}
  \vspace{2cm}
  {\large Program Studi Teknik Informatika\par}
  {\large Fakultas Teknologi Informasi\par}
  {\large Universitas Bale Bandung\par}
  \vfill
  {\large 2026\par}
\end{titlepage}

\cleardoublepage

% ============================================================
% Kata Pengantar
% ============================================================
\chapter*{Kata Pengantar}
\addcontentsline{toc}{chapter}{Kata Pengantar}

Puji syukur kehadirat Tuhan Yang Maha Esa atas terselesaikannya buku ajar \textit{Pemrograman Web} ini. Buku ini disusun dengan pendekatan \textbf{Outcome-Based Education (OBE)}, yang berfokus pada pencapaian kompetensi terukur sesuai dengan Capaian Pembelajaran Lulusan (CPL) dan Capaian Pembelajaran Mata Kuliah (CPMK).

Pemrograman Web merupakan keterampilan fundamental yang harus dikuasai oleh setiap mahasiswa Teknik Informatika. Mata kuliah ini tidak hanya mengajarkan cara membuat halaman web, tetapi juga cara berpikir dalam merancang dan mengimplementasikan aplikasi web yang meliputi analisis kebutuhan, wireframe dan mockup (UI/UX), antarmuka (HTML, CSS, JavaScript), logika server-side, basis data, integrasi API, pengujian, keamanan, dan deployment.

Buku ini dirancang untuk mendukung pembelajaran mahasiswa dengan pendekatan student\-centered learning. Setiap bab dilengkapi dengan:
\begin{itemize}
  \item Sub-CPMK yang jelas dan terukur
  \item Materi pokok dengan contoh kode HTML, CSS, JavaScript, dan PHP
  \item Aktivitas pembelajaran yang mendorong eksplorasi mandiri
  \item Latihan dan refleksi untuk penguatan pemahaman
  \item Asesmen untuk mengukur pencapaian kompetensi
  \item Checklist kompetensi untuk self-assessment
\end{itemize}

Kami berharap buku ini dapat menjadi panduan yang efektif dalam perjalanan pembelajaran Anda menguasai Pemrograman Web.

\vspace{1cm}
\begin{flushright}
Penyusun\\
\lbrack Nama Dosen, Gelar\rbrack
\end{flushright}

\cleardoublepage

% ============================================================
% Cara Menggunakan Buku Ini
% ============================================================
\chapter*{Cara Menggunakan Buku Ini}
\addcontentsline{toc}{chapter}{Cara Menggunakan Buku Ini}

Buku ajar ini dirancang dengan pendekatan OBE untuk memaksimalkan pencapaian pembelajaran Anda. Berikut panduan penggunaan buku ini:

\section*{Struktur Buku}

\textbf{Bab I: Pendahuluan dan Orientasi}\\
Memperkenalkan tujuan buku, keterkaitan dengan RPS, dan konteks kurikulum OBE.

\textbf{Bab II: Landasan Teori dan Konsep Dasar}\\
Menyajikan fondasi arsitektur web dan protokol HTTP yang menjadi basis pembelajaran.

\textbf{Bab III: Analisis Kebutuhan dan Perancangan Antarmuka}\\
Kebutuhan fungsional dan non-fungsional, wireframe, mockup, aksesibilitas, dan UX.

\textbf{Bab IV--V: HTML5}\\
Struktur, elemen semantik, link, gambar, tabel, form, dan media.

\textbf{Bab VI--VII: CSS3}\\
Selector, box model, layout dasar, Flexbox, Grid, dan desain responsif.

\textbf{Bab VIII--IX: JavaScript}\\
Sintaks, DOM, event handling, validasi form, dan library/framework front-end.

\textbf{Bab X--XII: Back-End dan Integrasi}\\
PHP (pemrosesan form, OOP), MySQL (CRUD, koneksi PHP), integrasi API dan koneksi front--back-end.

\textbf{Bab XIII--XV: Pengujian, Keamanan, dan Deployment}\\
Pengujian dan evaluasi performa, keamanan web, deployment dan administrasi dasar.

\textbf{Bab XVI: Evaluasi, Refleksi, dan Integrasi Kompetensi}\\
Berisi asesmen komprehensif dan panduan refleksi untuk mengukur pencapaian kompetensi.

\textbf{Lampiran}\\
Menyediakan rubrik penilaian, glosarium istilah pemrograman web, dan referensi tambahan.

\section*{Komponen dalam Setiap Bab}

\begin{enumerate}
  \item \textbf{Sub-CPMK}: Baca dengan seksama untuk memahami kompetensi yang harus dicapai
  \item \textbf{Materi Pokok}: Pelajari dengan cermat, jalankan semua contoh kode
  \item \textbf{Aktivitas Pembelajaran}: Lakukan secara mandiri atau berkelompok
  \item \textbf{Latihan}: Kerjakan untuk menguji pemahaman Anda
  \item \textbf{Asesmen}: Gunakan untuk mengukur pencapaian Sub-CPMK
  \item \textbf{Checklist}: Centang setelah yakin menguasai setiap indikator
\end{enumerate}

\section*{Tips Belajar Efektif}

\begin{itemize}
  \item Gunakan browser dan VS Code (dengan addon HTML/CSS/JS/PHP) untuk praktik \cite{vscode}
  \item Siapkan XAMPP untuk menjalankan PHP dan MySQL secara lokal \cite{xampp}
  \item Manfaatkan alat asisten coding (mis. GitHub Copilot) sesuai etika akademik bila tersedia \cite{vscode}
  \item Jalankan contoh kode di setiap bab dan modifikasi untuk eksperimen
  \item Manfaatkan dokumentasi MDN, W3Schools, dan sumber daring di daftar referensi
  \item Kerjakan proyek mini untuk mengintegrasikan HTML, CSS, JavaScript, dan server-side
  \item Diskusikan konsep yang sulit dengan teman atau dosen
  \item Manfaatkan checklist untuk self-assessment berkala
\end{itemize}

\cleardoublepage

% ============================================================
% Identitas Mata Kuliah
% ============================================================
\chapter*{Identitas Mata Kuliah}
\addcontentsline{toc}{chapter}{Identitas Mata Kuliah}

\begin{tabular}{ll}
  Nama Program Studi & : Teknik Informatika \\
  Nama Mata Kuliah & : Pemrograman Web \\
  Kode Mata Kuliah & : \lbrack Diisi oleh Program Studi\rbrack \\
  Semester & : \lbrack Diisi sesuai kurikulum\rbrack \\
  SKS / Bobot Kredit & : 2 SKS \\
  Jumlah Pertemuan & : 16 pertemuan \\
  Komposisi & : 70\% Praktek, 30\% Teori \\
  Dosen Pengampu & : \lbrack Nama Dosen, Gelar\rbrack \\
  Tanggal Penyusunan & : \lbrack Tanggal\rbrack \\
\end{tabular}

\vspace{1cm}

\section*{Capaian Pembelajaran Lulusan (CPL)}

CPL yang dibebankan pada mata kuliah ini:

\begin{enumerate}
  \item \textbf{CPL04:} Kemampuan mendesain, mengimplementasi dan mengevaluasi solusi berbasis computing yang memenuhi kebutuhan-kebutuhan computing pada sebuah disiplin program.
  
  \item \textbf{CPL05:} Mengidentifikasi dan menganalisis kebutuhan-kebutuhan pengguna dan mempertimbangkannya dalam memilih, membuat, mengintegrasi, mengevaluasi, dan mengadministrasi sistem berbasis computing.
\end{enumerate}

\section*{Capaian Pembelajaran Mata Kuliah (CPMK)}

Kemampuan atau kompetensi spesifik yang diharapkan mahasiswa kuasai setelah menyelesaikan mata kuliah:

\begin{enumerate}
  \item \textbf{CPMK-1:} Mahasiswa mampu mengidentifikasi kebutuhan fungsional dan non-fungsional pengguna untuk merancang solusi web yang intuitif melalui pembuatan wireframe dan mockup (UI/UX) yang mempertimbangkan aspek aksesibilitas dan pengalaman pengguna.
  
  \item \textbf{CPMK-2:} Mahasiswa mampu mendesain dan mengimplementasikan antarmuka web yang responsif dan dinamis menggunakan teknologi standar (HTML5, CSS3, JavaScript) serta framework modern untuk memenuhi spesifikasi perangkat pengguna yang beragam.
  
  \item \textbf{CPMK-3:} Mahasiswa mampu membangun arsitektur sisi server (back-end), mengelola basis data, dan mengintegrasikan API untuk menciptakan sistem berbasis computing yang fungsional dan aman sesuai dengan disiplin program.
  
  \item \textbf{CPMK-4:} Mahasiswa mampu melakukan pengujian (testing), mengevaluasi keamanan serta performa aplikasi web, dan melakukan proses deployment ke server produksi (administrasi sistem) untuk memastikan aplikasi dapat diakses secara optimal oleh publik.
\end{enumerate}

\section*{Matriks Keterkaitan CPL dan CPMK}

\begin{table}[htbp]
\centering
\small
\begin{tabular}{|c|c|c|>{\raggedright\arraybackslash}p{4.5cm}|}
  \hline
  \textbf{CPL} & \textbf{CPMK} & \textbf{Kontribusi} & \textbf{Keterangan} \\
  \hline
  CPL04 & CPMK-1 & Tinggi & Analisis kebutuhan dan perancangan antarmuka \\
  \hline
  CPL04 & CPMK-2 & Tinggi & Implementasi front-end (HTML5, CSS3, JavaScript, framework) \\
  \hline
  CPL04 & CPMK-3 & Tinggi & Back-end, basis data, integrasi API \\
  \hline
  CPL04 & CPMK-4 & Tinggi & Pengujian, keamanan, deployment \\
  \hline
  CPL05 & CPMK-1 & Tinggi & Identifikasi dan analisis kebutuhan pengguna \\
  \hline
  CPL05 & CPMK-2, CPMK-3, CPMK-4 & Tinggi & Integrasi, evaluasi, dan administrasi sistem web \\
  \hline
\end{tabular}
\caption{Matriks Keterkaitan CPL dan CPMK}
\end{table}

\cleardoublepage

% ============================================================
% Daftar Isi
% ============================================================
\phantomsection
\addcontentsline{toc}{chapter}{Daftar Isi}
\tableofcontents
\cleardoublepage

\clearpage
\phantomsection
\addcontentsline{toc}{chapter}{Daftar Gambar}
\listoffigures
\cleardoublepage

\phantomsection
\addcontentsline{toc}{chapter}{Daftar Tabel}
\listoftables
\cleardoublepage

\phantomsection
\addcontentsline{toc}{chapter}{Daftar Kode Program}
\lstlistoflistings
\cleardoublepage
